\documentclass[12pt,a4paper]{article}
\usepackage[utf8]{vietnam}
\usepackage{amsmath,amssymb}
\usepackage{geometry}
\geometry{
  top=12mm,
  bottom=12mm,
  left=14mm,
  right=14mm,
  footskip=6mm,
  headheight=14pt
}

\usepackage{tikz}
\usetikzlibrary{arrows.meta,calc,shapes.geometric}
\usepackage[most]{tcolorbox}
\usepackage{fontawesome}
\usepackage{fancyhdr}
\usepackage{enumitem}

\pagestyle{fancy}
\fancyhf{}
\fancyhead[L]{\footnotesize\color{teal!80!black}\faBook\ \textbf{SỔ TAY KIẾN THỨC TOÁN 3}}
\fancyhead[R]{\footnotesize\color{teal!80!black}\textbf{Lớp Toán Thầy Lĩnh -- loptoanthaylinh.vn}}
\fancyfoot[C]{\footnotesize\color{gray}-- Sổ tay kiến thức: Bài 14. Một phần mấy (1/2, 1/3, 1/4, 1/5...) --}
\renewcommand{\headrulewidth}{0.4pt}
\renewcommand{\footrulewidth}{0pt}

\setlength{\parindent}{0pt}
\setlength{\parskip}{0pt}

\begin{document}

% --- BANNER BÀI HỌC ---
\begin{tcolorbox}[
  enhanced,
  colback=teal!8!white,
  colframe=teal!70!black,
  arc=3mm, boxrule=1.2pt,
  left=3mm, right=3mm, top=2mm, bottom=2mm,
  halign=center
]
  {\LARGE\bfseries\color{teal!90!black}\faBookmark\ BÀI 14: MỘT PHẦN MẤY (1/2, 1/3, 1/4, 1/5...)}\\[1mm]
  {\itshape\footnotesize\color{black!75} Bản chất chia đều các phần bằng nhau, nhận biết và đọc viết một phần mấy trên 1 trang A4}
\end{tcolorbox}

\vspace{-2mm}

\noindent
\begin{minipage}[t]{0.49\textwidth}\vspace{0pt}
\begin{tcolorbox}[
  enhanced,
  colback=blue!4!white,
  colframe=blue!75!black,
  arc=2mm, boxrule=0.8pt,
  left=2mm, right=2mm, top=1.8mm, bottom=1.8mm,
  title={\footnotesize\bfseries\color{white}\faCheckSquare\ 1. Khái niệm 'Một phần mấy'}
]
{\footnotesize
\begin{itemize}\setlength{\itemsep}{1pt}\setlength{\leftskip}{-2.5mm}
  \item \textbf{Bản chất:} Chia một hình (hoặc một nhóm đồ vật) thành \textbf{$n$ phần bằng nhau}. Lấy đi $1$ phần, ta được \textbf{một phần $n$}, viết là $\mathbf{\dfrac{1}{n}}$.
  \item \textbf{Cấu tạo kí hiệu $\dfrac{1}{n}$:}
  \begin{itemize}\setlength{\itemsep}{0.5pt}\setlength{\leftskip}{-2.5mm}
    \item Số $1$ ở trên: Số phần được lấy ra.
    \item Vạch ngang ở giữa: Thể hiện phép chia đều.
    \item Số $n$ ở dưới: Tổng số phần bằng nhau được chia.
  \end{itemize}
  \item $\dfrac{1}{2}$ còn gọi là \textbf{một nửa}.
\end{itemize}
}
\end{tcolorbox}
\end{minipage}
\hfill
\begin{minipage}[t]{0.49\textwidth}\vspace{0pt}
\begin{tcolorbox}[
  enhanced,
  colback=teal!4!white,
  colframe=teal!75!black,
  arc=2mm, boxrule=0.8pt,
  left=2mm, right=2mm, top=1.8mm, bottom=1.8mm,
  title={\footnotesize\bfseries\color{white}\faCalculator\ 2. Kĩ thuật tìm $\dfrac{1}{n}$ của một số}
]
{\footnotesize
\begin{itemize}\setlength{\itemsep}{1pt}\setlength{\leftskip}{-2.5mm}
  \item \textbf{Quy tắc vàng:} Muốn tìm một phần $n$ của số $A$, ta \textbf{lấy số đó chia cho $n$}:
  $$\text{Giá trị cần tìm} = A : n$$
  \item \textit{Ví dụ:}
  \begin{itemize}\setlength{\itemsep}{0.5pt}\setlength{\leftskip}{-2.5mm}
    \item Tìm $\dfrac{1}{3}$ của $18\text{ cm}$: Lấy $18 : 3 = \mathbf{6\text{ cm}}$.
    \item Tìm $\dfrac{1}{5}$ của $35\text{ kg}$: Lấy $35 : 5 = \mathbf{7\text{ kg}}$.
    \item Tìm $\dfrac{1}{4}$ của $24$ giờ: Lấy $24 : 4 = \mathbf{6}$ giờ.
  \end{itemize}
\end{itemize}
}
\end{tcolorbox}
\end{minipage}

\vspace{-2.5mm}
% --- KHỐI 3: SƠ ĐỒ TƯ DUY TIKZ ---
\noindent
\begin{tcolorbox}[
  enhanced,
  colback=orange!4!white,
  colframe=orange!80!black,
  arc=2mm, boxrule=0.8pt,
  left=2mm, right=2mm, top=1.5mm, bottom=1.5mm,
  title={\footnotesize\bfseries\color{white}\faSitemap\ SƠ ĐỒ TƯ DUY: MÔ HÌNH TRỰC QUAN MỘT PHẦN MẤY}
]
\centering
\begin{tikzpicture}[
  part/.style={draw=blue!80!black, minimum height=0.55cm, font=\scriptsize\bfseries}
]
  % 1/2
  \node[anchor=east, font=\scriptsize\bfseries] at (0, 0.9) {$\dfrac{1}{2}$ (Chia 2 phần):};
  \node[part, fill=teal!30, minimum width=1.8cm] at (1.2, 0.9) {$\dfrac{1}{2}$};
  \node[part, fill=white, minimum width=1.8cm] at (3.0, 0.9) {};

  % 1/3
  \node[anchor=east, font=\scriptsize\bfseries] at (0, 0.3) {$\dfrac{1}{3}$ (Chia 3 phần):};
  \node[part, fill=teal!30, minimum width=1.2cm] at (0.9, 0.3) {$\dfrac{1}{3}$};
  \node[part, fill=white, minimum width=1.2cm] at (2.1, 0.3) {};
  \node[part, fill=white, minimum width=1.2cm] at (3.3, 0.3) {};

  % 1/4
  \node[anchor=east, font=\scriptsize\bfseries] at (0, -0.3) {$\dfrac{1}{4}$ (Chia 4 phần):};
  \node[part, fill=teal!30, minimum width=0.9cm] at (0.75, -0.3) {$\dfrac{1}{4}$};
  \node[part, fill=white, minimum width=0.9cm] at (1.65, -0.3) {};
  \node[part, fill=white, minimum width=0.9cm] at (2.55, -0.3) {};
  \node[part, fill=white, minimum width=0.9cm] at (3.45, -0.3) {};
\end{tikzpicture}
\end{tcolorbox}

\vspace{-2.5mm}
% --- KHỐI 4: VÍ DỤ & MẸO TRÁNH BẪY ---
\noindent
\begin{minipage}[t]{0.49\textwidth}\vspace{0pt}
\begin{tcolorbox}[
  enhanced,
  colback=green!4!white,
  colframe=green!60!black,
  arc=2mm, boxrule=0.8pt,
  left=2mm, right=2mm, top=1.8mm, bottom=1.8mm,
  title={\footnotesize\bfseries\color{white}\faCheckCircle\ Ví dụ thực hành chuẩn mực}
]
{\footnotesize
\textbf{Bài toán:}\\
Lớp 3A có $32$ học sinh, trong đó $\dfrac{1}{4}$ số học sinh tham gia câu lạc bộ Cờ vua. Hỏi có bao nhiêu bạn tham gia câu lạc bộ Cờ vua?\\
\textbf{Bài giải:}\\
Số học sinh tham gia câu lạc bộ Cờ vua là:
$$32 : 4 = \mathbf{8}\text{ (học sinh)}$$
\textbf{Đáp số:} $8$ học sinh.
}
\end{tcolorbox}
\end{minipage}
\hfill
\begin{minipage}[t]{0.49\textwidth}\vspace{0pt}
\begin{tcolorbox}[
  enhanced,
  colback=red!4!white,
  colframe=red!70!black,
  arc=2mm, boxrule=0.8pt,
  left=2mm, right=2mm, top=1.8mm, bottom=1.8mm,
  title={\footnotesize\bfseries\color{white}\faExclamationTriangle\ Mẹo tránh bẫy trắc nghiệm}
]
{\footnotesize
\begin{itemize}\setlength{\itemsep}{1.5pt}\setlength{\leftskip}{-2.5mm}
  \item \textbf{Bẫy chia không đều:}\\
  Muốn xác định hình đã tô màu $\dfrac{1}{3}$, các phần phải \textbf{hoàn toàn bằng nhau}. Nếu chia 3 phần to nhỏ khác nhau thì KHÔNG phải $\dfrac{1}{3}$.
  \item \textbf{Bẫy nhầm phép tính:}\\
  Tìm $\dfrac{1}{5}$ của $30$ là làm phép \textbf{chia} $30 : 5 = 6$, nhiều bạn nhầm sang phép trừ $30 - 5 = 25$!
\end{itemize}
}
\end{tcolorbox}
\end{minipage}


\end{document}
